This chapter should clearly motivate the need for gravitational wave measurements in the context of cosmology and astrophysics. It will (briefly) explain the types of sources that the detectors will be measuring (including a bit about LISA). It will talk about tests of general relativity, measurements of the age of the universe, multi-messenger measurements of the neutron star equation of state (its a combination of basically every field of physics), and so on... (I will read up more specific examples, but having large numbers of observations will tell us something about astrophysics). Although I'm not sure if I will mention dark matter in great depth, it probably deserves a mention since its one of the great mysteries of physics and gravitational wave detections may give insights into it. This will be an expanded version of the introduction to most papers written about gravitational wave detections, which hopefully will give you an idea of the tone of this chapter. It will be a relatively short chapter (about 5-10 pages?), but its vital in giving the rest of the work a reason to exist (some people do study graviational wave detectors as tools for probing the quantum nature of light, but most people will not see this as a good enough reason to spend billions of dollars developing them, heck not even everyone in the LVK collaboration thinks this is a good motivation for building the instruments!). Finally it will end with an explanation of each chapter in the thesis. This chapter will not cover the working principles of an interferometer
%\section{Introduction}
%\section{general relativity}
%\section{What is a gravitational wave}
%\section{How detections of gravitational waves have verified GR}
%\section{Sources of gravitational waves}
%formula for neutron star distance
%\section{how gravitational wave detections have been important to cosmology}
%\section{how gravitational wave detections have been important to astronomy}
%\section{thesis structure}
